\section{Representations of Lie Groups and Lie Algebras}

  Let $\mathcal{G}$ be a Lie group and let
  \begin{equation}
    \rho: \mathcal{G} \rightarrow \GL(V)
  \end{equation}
  be the representation of $\mathcal{G}$. Then, let $\mathfrak{g}$ be the Lie algebra of $\mathcal{G}$, and $\mathfrak{gl}(V)$ be the Lie algebra of $\GL(V)$. Then, $\rho$ induces another homomorphism 
  \begin{equation}
    \varrho: \mathfrak{g} \rightarrow \mathfrak{gl}(V)
  \end{equation}
  where the bracket structure (in this case, the comutator in the matrix algebra) is preserved. 
  \begin{equation}
    \varrho([X,Y]) = [\varrho(X), \varrho(Y)]
  \end{equation}
  We can visualize this induced homomorphism with the following commutative diagram, which states that $\rho \circ \exp = \exp \circ \varrho$. 

  \begin{equation}
    \begin{tikzcd}
      \mathcal{G} \arrow{r}{\rho} & \GL(V)\\
      \mathfrak{g} \arrow{u}{\exp} \arrow{r}{\varrho} & \mathfrak{gl}(V) \arrow{u}{\exp}
      \end{tikzcd}
  \end{equation}

  Note that there are very crucial differences between $\rho$ and $\varrho$. First, $\rho$ is a homomorphism between \textbf{groups}, while $\varrho$ is a homomorphism between \textbf{\hyperref[lin-def:vector-space]{vector spaces}}. Additionally, $\GL(V)$ is a group, not a linear space, while $\mathfrak{gl}(V)$ is a linear space. Finally, note that $\GL(V)$ is restricted to only matrices with nonzero determinants, while the elements of $\mathfrak{gl}(V)$ can be any matrix. 

  \begin{example}[Matrix Representation of $SE(n)$]
    \label{ex:matrix-representation-of-sen}
    The representation of SE$(n)$ to $\GL(n+1, \mathbb{R})$ and $\mathfrak{se}(n)$ to $\mathfrak{gl}_{n+1} \mathbb{R}$ induces the second homomorphism $\varrho: \mathfrak{gl}_{n+1} \mathbb{R} \rightarrow \GL(n+1, \mathbb{R})$. 
  \end{example}

  For our case, we will consider that any representation can be written as a direct sum of \hyperref[def:irreducible-representation]{irreducible representations}. We will now proceed to find an irreducible representation of $\mathfrak{sl}_2 \mathbb{C}$. This means that we want to find the smallest (lowest dimensional) \hyperref[lin-def:vector-space]{vector space} $V$ such that there exists a representation
  \begin{equation}
    \varrho: \mathfrak{sl}_2 \mathbb{C} \rightarrow \mathfrak{gl}(V)
  \end{equation}
  We will write, as shorthand notation, that 
  \begin{equation}
    H = \varrho(H), X = \varrho(X), Y = \varrho(Y)
  \end{equation}
  Clearly, $H, X, Y \in \mathfrak{gl}(V) \simeq \mathfrak{gl}(\mathbb{C}^n)$. By the spectral theorem, we can find an orthonormal basis of eigenvectors $e_1, e_2, ..., e_n$ of the mapping $H$ such that
  \begin{equation}
    H e_i = \lambda_i e_i, \; \lambda_i \in \mathbb{C}
  \end{equation}
  Since $[H,X] = 2X$, it follows that $HX e_i - X H e_i = 2X e_i \implies H (X e_i) = (\lambda_i + 2) (X e_i) \implies Xe_i$ for all $i = 1, 2, ..., n$ are also eigenvectors of $H$ with eigenvalue $(\lambda_i + 2)$, or $X e_i = 0$. So, $X$ is a "ladder operator" that maps each eigenvector $e_i$ with eigenvalue $\lambda_i$ to a different eigenvector $e_j$ with eigenvalue $\lambda_j = \lambda_i + 2$. Having nowhere to be mapped to, the eigenvector with the largest eigenvalue (which must exist since $V$ is finite dimensional) will get mapped to the $0$ vector by $X$. Let us denote this eigenvector having the maximum eigenvalue $m$, as $v_m$. 

  Similarly, $[H,Y] = -2Y$ implies that
  \begin{equation}
    HY e_i - YH e_i = -2Y e_i \implies H(Y e_i) = (\lambda_i - 2)(Y e_i)
  \end{equation}

  implying that $Y$ maps each eigenvector $e_i$ with eigenvaue $\lambda_i$ to another eigenvector $e_j$ with eigenvalue $\lambda_j = \lambda_i - 2$, except for the eigenvector with smallest eigenvalue, which gets mapped to $0$. Since $Y$ clearly maps each eigenvector to a different eigenvector that has a strictly decreasing eigenvalue, we can construct a basis of $V$ to be
  \begin{equation}
    \{v_m, Y v_m, Y^2 v_m, Y^3 v_m, ..., Y^{n-1} v_m\}
  \end{equation}
  (remember that $Y^n v_m = 0$). So, elements of $\mathfrak{sl}_2 \mathbb{C}$ acts on the space $V$ with basis above. To continue, we introduce the following theorem. 

  \begin{theorem}[Action of the Raising and Lowering Operators]
    \label{thm:action-of-the-raising-and-lowering-operators}
    \begin{equation}
      X Y^j v_m = j(m-j+1) Y^{j-1} v_m
    \end{equation}
  \end{theorem}
  \begin{proof}
    By induction on $j$ using bracket relations.
  \end{proof}

  $V$ is $n$-dimensional. Since $Y^n v_m = 0$ and $Y^{n-1} v_m \neq 0$, we use the theorem above to get
  \begin{equation}
    0 = X Y^n v_m = n (m-n+1) Y^{n-1} v_m \implies m-n+1=0
  \end{equation}
  So, $n = m+1$, which means that the eigenvalues of $H$ are
  \begin{equation}
    m, m-2, m-4, \ldots, m - 2(n-1) = -m
  \end{equation}

  and we are done. We now classify the 1, 2, and 3 dimensional \hyperref[def:irreducible-representation]{irreducible representations} of $\mathfrak{sl}_2 \mathbb{C}$. 
  \begin{enumerate}
    \item When $n = 1$ (i.e. dimension is 1), $m = n-1 = 0$, meaning that the greatest (and only) eigenvalue is $0$. That is, 
      \begin{equation}
        H v_0 = 0,\; X v_0 = 0,\; Y v_0 = 0
      \end{equation}
    which is the trivial representation of $\mathfrak{sl}_2 \mathbb{C}$. Explicitly, we can completely define the representation (which is a linear homomorphism) with the three equations. 
    \begin{equation}
      \varrho(H) = (0),\; \varrho(X) = (0),\; \varrho(Y) = (0)
    \end{equation}

    \item When $n = 2$ and $m=1$. We now look for a 2 dimensional irreducible representation. The eigenvalues are $1$ and $-1$, with $\{v_1, v_{-1}\}$ as a basis of 2 dimensional space $V$. Then we have 
      \begin{align*}
        & Hv_1 = v_1, \; Hv_{-1} = - v_{-1} \\
        & X v_1 = 0, \; X v_{-1} = v_1 \\
        & Y v_1 = v_{-1}, \; Y v_{-1} = 0
      \end{align*}
    which explicitly translates to the representation $\varrho$ being defined
    \begin{equation}
      \varrho(H) = \begin{pmatrix}
      1&0\\0&-1
      \end{pmatrix}, \; \begin{pmatrix}
      0&1\\0&0
      \end{pmatrix}, \; \begin{pmatrix}
      0&0\\1&0
      \end{pmatrix}
    \end{equation}

    \item When $n=3 \implies m=2$, the basis is $\{v_{-2}, v_0, v_2\}$ with eigenvalues $2, 0, -2$, and the irreducible representation $\varrho$ is defined
      \begin{equation}
        \varrho(H) = \begin{pmatrix}
        2&&\\&0&\\&&-2
        \end{pmatrix}, \varrho(Y) = \begin{pmatrix}
        0&0&0\\1&0&0\\0&1&0
        \end{pmatrix}, \varrho(X) = \begin{pmatrix}
        0&1&0\\0&0&1\\0&0&0
        \end{pmatrix}
      \end{equation}

    \item The same process continues on for $n=4, 5, ...$, and this entirely classifies the irreducible representations of $\mathfrak{sl}_2 \mathbb{C}$. 
  \end{enumerate} 


  The adjoint maps $\mathrm{Ad}$ and $\mathrm{ad}$ defined in the \hyperref[dg-def:the-adjoint-map-of-a-lie-group]{notes on differential geometry} are also examples of representations: $\mathrm{Ad}$ is a representation of $\mathcal{G}$ on its Lie algebra $\mathfrak{g}$, and $\mathrm{ad}$ is the induced representation of $\mathfrak{g}$ on itself. 
